\chapter{I ritardi}
\label{chap:ritardi}

Oltre alle dissonanze "strutturali" contenute negli accordi di settima, l'armonia si arricchisce attraverso l'uso di dissonanze melodiche, ovvero note estranee all'accordo che creano una tensione momentanea prima di risolvere. Il \textbf{ritardo} è una delle figure più importanti ed espressive di questa categoria.

\section{Definizione e meccanismo del ritardo}
\label{sec:definizione_ritardo}

Il ritardo è una dissonanza che si manifesta sul tempo forte di una misura (o su una parte forte di un tempo), ritardando, appunto, l'arrivo di una nota reale dell'accordo. La sua efficacia espressiva deriva dal rispetto di un rigoroso processo temporale in tre fasi.

\begin{definizione}[Le tre fasi del Ritardo]
	\label{def:fasi_ritardo}
	Un ritardo si articola sempre in tre momenti distinti:
	\begin{enumerate}
		\item \textbf{Preparazione:} La nota che diventerà dissonante viene introdotta come nota consonante nell'accordo precedente, preferibilmente nella stessa voce e alla stessa altezza.
		\item \textbf{Percussione:} L'accordo cambia, ma la nota preparata viene tenuta (legata), diventando così dissonante rispetto al nuovo contesto armonico. Questa fase deve avvenire su un tempo forte.
		\item \textbf{Risoluzione:} La nota dissonante scende di grado congiunto su una nota reale dell'accordo.
	\end{enumerate}
\end{definizione}

\section{Il ritardo 4-3}
\label{sec:ritardo_4_3}

Il ritardo 4-3 è il più celebre e utilizzato. Ritarda l'arrivo della terza dell'accordo: sulla dominante ritarda la sensibile (Si in Do maggiore), sulla tonica ritarda la terza dell'accordo (Mi in Do maggiore).

\subsection{Il ritardo nella cadenza perfetta}
\label{subsec:ritardo_cadenza}

Il ritardo 4-3 trova la sua applicazione più efficace nella cadenza perfetta (V-I), dove può amplificare la tensione della dominante e ritardare la risoluzione sulla tonica. Come per la settima di dominante, esistono due forme principali di realizzazione.

\begin{esempio}[Forma A: V7 completo → I incompleto con ritardo 4-3 in Do Maggiore]
	\label{ex:ritardo_completo_maggiore}
	Questa progressione utilizza tutti e quattro i suoni sulla dominante (compresa la settima) e risolve su un accordo di tonica incompleto (senza quinta), con raddoppio della fondamentale e ritardo della terza.
	
	\musica{es_cap10_01}
	
	La preparazione del ritardo avviene nell'accordo di IV: il Fa (nel contralto) è nota reale dell'accordo. Con l'arrivo del V7 (Sol al basso) il Fa è tenuto (legato) e diventa la settima; il soprano scende alla sensibile Si e il tenore salta al Re (quinta). Nella seconda battuta, sul basso Do, il Fa è percosso come ritardo 4 e risolve su Mi (3), mentre la sensibile sale a Do e il Re scende a Do.
\end{esempio}

\begin{esempio}[Forma B: V7 incompleto → I completo con ritardo 4-3 in Do Maggiore]
	\label{ex:ritardo_incompleto_maggiore}
	Questa progressione omette la quinta dalla dominante, permettendo di risolvere su un accordo di tonica completo con ritardo della terza.
	
	\musica{es_cap10_02}
	
	In questa versione la dominante è incompleta (senza Re, la quinta) e la fondamentale è raddoppiata; la settima (Fa) è tenuta dal contralto e diventa il ritardo. L'accordo di tonica risultante è completo, con fondamentale, terza (ritardata) e quinta.
\end{esempio}

\subsection{Il ritardo nel modo minore}
\label{subsec:ritardo_minore}

Nel modo minore, il ritardo 4-3 segue gli stessi principi, ma si adatta alle caratteristiche della scala minore armonica, in particolare all'uso della sensibile innalzata nel V7.

\begin{esempio}[Forma A: V7 completo → I incompleto con ritardo 4-3 in La minore]
	\label{ex:ritardo_completo_minore}
	Progressione con V7 completo che risolve su I incompleto con ritardo.
	
	\musica{es_cap10_03}
\end{esempio}

\begin{esempio}[Forma B: V7 incompleto → I completo con ritardo 4-3 in La minore]
	\label{ex:ritardo_incompleto_minore}
	Progressione con V7 incompleto che risolve su I completo con ritardo.
	
	\musica{es_cap10_04}
\end{esempio}

\begin{regola}[Cifratura del Ritardo]
	\label{reg:cifratura_ritardo}
	Nel basso figurato, il ritardo 4-3 si indica scrivendo in sequenza \texttt{4 3}. Il numero 4 rappresenta la dissonanza (la quarta rispetto al basso), mentre il 3 indica la sua risoluzione (la terza dell'accordo).
\end{regola}

Gli esempi mostrano come il ritardo 4-3 possa arricchire la progressione cadenzale, creando uno slancio melodico nelle voci superiori mentre il basso mantiene la sua funzione armonica di sostegno. La scelta tra forma completa e incompleta dipende dal contesto musicale e dalla necessità di preservare o meno tutte le note dell'accordo di arrivo.